
\chapter{Considerações finais}

\setlength{\epigraphwidth}{0.9\textwidth}

\epigraph{A questão [da maquinaria] era central para as relações cotidianas entre patrão e trabalhador, mas também despertava grande interesse teórico e ideológico. A própria tecnologia que sustentava a economia e a sociedade constituía um terreno de desafios e disputas. A máquina foi amplamente debatida em todos os setores da sociedade. Ela provocava tanto o clérigo da aldeia quanto o intelectual cosmopolita; preocupava tanto o político quanto o trabalhador e o empregador; tanto o reformador social [\textit{social reformer}] quanto o cientista e o inventor. Esses grupos debatiam os custos e benefícios da nova tecnologia. Saudavam a libertação que ela proporcionava em relação aos limites ao crescimento, mas discordavam quanto ao impacto que teria sobre salários, emprego e qualificação profissional. Especulavam sobre as mudanças que a máquina traria às relações sociais — mudanças que, posteriormente, ora acolhiam com entusiasmo, ora temiam. Até mesmo as origens e a propriedade do maquinário foram postas em questão. Havia entusiasmo e temor diante dessa força desconhecida que avançava implacavelmente, deixando para trás a antiga sociedade.}{\parencite[Berg \textit{apud}][p. 80-1]{Pasquinelli2023}}